% TOP_PROBLEM: {"id": 4200002, "title": "N-body problem", "category_id": 11, "category": {"id": 11, "name": "dynamical_systems", "display_name": "Dynamical Systems", "description": "Problems about long-term behavior of deterministic systems, Hamiltonian dynamics, and periodic orbits.", "slug": "dynamical-systems", "order_index": 11, "created_at": "2026-07-31T15:26:25.671Z"}, "status": "partially_solved"}
% FIRST_POSTED: null
% ENTRY_KIND: statement_only
% Imported from: batch14.tex; UnsolvedMath ID 4200002
% Compile twice: lualatex 4200002.tex
\documentclass[11pt]{article}
\usepackage{fontspec}
\setmainfont{Latin Modern Roman}
\usepackage{amsmath,amssymb,mathtools,unicode-math}
\setmathfont{Latin Modern Math}
\usepackage{xurl,hyperref}
\hypersetup{hidelinks,pdftitle={UnsolvedMath Problem 4200002}}
\newcommand{\sref}[2]{\hyperlink{source:#1:#2}{[S#2]}}
\newcommand{\cataloguescope}[1]{\paragraph{Mathematical question.}#1}
\begin{document}
% UnsolvedMath ID 4200002; source catalogue code AMR-041-0002
\setcounter{section}{4200001}
\section{N-body problem}\label{4200002}
{\small\sffamily UnsolvedMath ID 4200002 · Dynamical Systems}
\subsection{Definitions and mathematical statement}

\paragraph{Shared notation.}$\mathbb N=\{1,2,\ldots\}$, and $\mathbb Z,\mathbb Q,\mathbb R,\mathbb C$ denote the integers, rationals, reals and complex numbers. Any other conventions are specified locally.


Fix an integer \(N\ge2\), positive masses \(m_1,\ldots,m_N\), and a gravitational constant \(G>0\). For the Newtonian \(N\)-body problem in \(\mathbb R^3\), the collision-free phase space is
\[\mathcal P=\{(q,p)\in(\mathbb R^3)^N\times(\mathbb R^3)^N:q_i\ne q_j\ (i\ne j)\}.\]
Positions \(q_i\) and momenta \(p_i\) evolve by Hamilton's equations for
\[H(q,p)=\sum_{i=1}^N\frac{|p_i|^2}{2m_i}
-G\sum_{i<j}\frac{m_i m_j}{|q_i-q_j|},\qquad
\dot q_i=p_i/m_i,\quad \dot p_i=-\nabla_{q_i}H.\]
For \(z\in\mathcal P\), let \(I_z\) be the maximal time interval of its collision-free classical solution containing time zero. The set of global initial conditions is \(\mathcal G=\{z:I_z=\mathbb R\}\). At a finite endpoint of \(I_z\), a collision singularity has a collision configuration as the limiting positions; a noncollision singularity is a finite-time failure of continuation without such a collision endpoint. Measure means Lebesgue measure on the full phase space, not on a fixed energy level.
\paragraph{Statement recorded in the supplied source.}
Determine the measure of \(\mathcal P\setminus\mathcal G\). In particular, for arbitrary \(N\) and positive masses, does \(\mathcal G\) have full Lebesgue measure in \(\mathcal P\)? The complementary set separates into collision and noncollision singularities; the central measure-zero question is whether initial conditions leading to noncollision singularities have Lebesgue measure zero. By time reversal, the full-measure question is equivalent to asking for almost-everywhere existence for all positive times. \sref{4200002}{2}
\subsection{Short English statement}
Do almost all collision-free initial conditions of the Newtonian N-body problem generate solutions for all time? Determine the measure of the exceptional collision and noncollision singularities.
\subsection{Sources}
\begin{description}
\item[\hypertarget{source:4200002:1}{S1}] ULAM.AI and the UnsolvedMath Contributors, UnsolvedMath: A Curated Collection of Open Mathematics Problems, record 4200002 (AMR-041-0002). CC BY 4.0. \url{https://huggingface.co/datasets/ulamai/UnsolvedMath}.
\item[\hypertarget{source:4200002:2}{S2}] Jinxin Xue, Noncollision Singularities in a Planar Four-body Problem (2014; revised 2020), arXiv:1409.0048, Section 1.1, Motivations and perspectives, Conjecture 1, p. 3. \url{https://arxiv.org/pdf/1409.0048}.
\end{description}
\label{end:4200002}
\end{document}
